\chapter{Stato dell'Arte}
\label{ch:stato-arte}

In questo capitolo vengono analizzate le principali soluzioni esistenti nel panorama della logistica collaborativa, del monitoraggio IoT della catena del freddo e della rilevazione di anomalie su serie temporali. Per ciascuna soluzione viene fornita una descrizione sintetica, seguita da un'analisi critica che ne evidenzia i limiti rispetto alla proposta presentata in questo lavoro.

\section{Flock Freight -- FlockDirect\textsuperscript{\textregistered} Shared Truckload}

Flock Freight è un'azienda statunitense che ha introdotto il concetto di \textit{Shared Truckload} (STL), una modalità di spedizione che consente di consolidare carichi di clienti diversi su un unico camion mediante una tecnologia brevettata basata su intelligenza artificiale. L'algoritmo di Flock Freight analizza posizioni geografiche, calendari di spedizione e dimensioni dei carichi per individuare, tra miliardi di combinazioni possibili, gli abbinamenti più efficienti, instradando il veicolo su percorsi diretti e privi di soste in magazzini intermedi \cite{flockfreight2024}. A differenza del tradizionale LTL (\textit{Less Than Truckload}), in cui la merce viene scaricata e ricaricata più volte lungo il percorso, nello STL il carico non lascia mai il veicolo fino alla consegna, riducendo drasticamente il rischio di danneggiamento e i tempi di transito. Flock Freight dichiara consegne puntuali nel 97\% dei casi e un tasso di danneggiamento 5,4 volte inferiore rispetto all'LTL, con una riduzione delle emissioni di CO\textsubscript{2} fino al 40\%.

Nonostante la solidità della soluzione, il modello di Flock Freight presenta caratteristiche differenti rispetto all'approccio proposto in ChillChain. Dal punto di vista dell'ottimizzazione, Flock adotta un approccio flessibile basato su AI che può combinare spedizioni con percorsi anche molto diversi tra loro, massimizzando il numero di abbinamenti possibili sull'intera rete. ChillChain segue invece un approccio spazialmente più rigido: le spedizioni condivise devono trovarsi lungo il percorso già programmato del veicolo, con una tolleranza massima di deviazione sia in termini di distanza dalla rotta sia di incremento della durata del viaggio. Questo vincolo riduce il numero di combinazioni possibili, ma quando il matching avviene il risparmio è potenzialmente più elevato, poiché la deviazione introdotta è minima e la riduzione di prezzo scala in modo aggressivo con il numero di spedizioni a bordo. Si tratta di due approcci alternativi, ciascuno con i propri compromessi tra flessibilità e ottimizzazione del costo. Inoltre, la piattaforma Flock Freight è progettata per il trasporto merci generico e non integra alcuna funzionalità di monitoraggio della catena del freddo: non offre strumenti di telemetria né di rilevazione anomalie per spedizioni refrigerate.

\section{Quicargo -- Piattaforma collaborativa con pricing dinamico}

Quicargo è una piattaforma logistica digitale fondata nei Paesi Bassi che opera come intermediario tra mittenti e trasportatori nel segmento LTL. 
La piattaforma consente ai clienti di prenotare spedizioni di pallet verso 29 Paesi europei \cite{quicargo2024}. Un aspetto particolarmente rilevante dal punto di vista scientifico è il lavoro di Atasoy, Schulte e Steenkamp (2020), che hanno studiato il caso Quicargo formulando un problema collaborativo di \textit{pickup and delivery} con finestre temporali (PDPTW), in cui il pricing pagato dalla piattaforma ai trasportatori viene modulato dinamicamente per bilanciare gli interessi di tutti gli attori coinvolti \cite{atasoy2020}. 

Quicargo opera esclusivamente nel trasporto merci generico. Il trasporto di prodotti refrigerati non è disponibile. ChillChain offre invece un servizio dedicato al trasporto refrigerato.

\section{Gillespie et al. -- Rilevazione anomalie IoT nel trasporto refrigerato}

Gillespie et al. (2023) propongono un sistema IoT per la rilevazione in tempo reale di anomalie durante il trasporto refrigerato di prodotti alimentari deperibili, pubblicato sulla rivista \textit{Sustainability} \cite{gillespie2023}. 
Il sistema si basa sul data logger cellulare Digital Matter Eagle, collegato a due sonde di temperatura posizionate all'interno del vano refrigerato del veicolo. I dati vengono trasmessi a una piattaforma cloud, dove un algoritmo di alerting basato su soglie di temperatura genera notifiche destinate al personale di magazzino e ai conducenti. 

Il lavoro di Gillespie et al. rappresenta un riferimento importante per il monitoraggio IoT applicato alla catena del freddo, in quanto dimostra la fattibilità di un sistema di alerting in tempo reale a basso costo basato su hardware commerciale. Tuttavia, la soluzione presenta un limite dal punto di vista della rilevazione delle anomalie: il sistema si basa esclusivamente su soglie statiche di temperatura. Questo approccio, pur essendo semplice e robusto, non è in grado di identificare pattern anomali complessi. ChillChain propone l'utilizzo di un modello LSTM Autoencoder per la rilevazione di anomalie multivariate per la cattura di anomalie complesse.

\section{Sgueglia et al. -- Systematic literature review su anomaly detection IoT}

Sgueglia, Di Sorbo, Visaggio e Canfora (2022) presentano una revisione sistematica della letteratura sulle soluzioni di rilevazione anomalie per serie temporali in ambienti IoT, pubblicata sulla rivista \textit{Future Generation Computer Systems} \cite{sgueglia2022}. Lo studio analizza 62 articoli pubblicati tra il 2014 e il 2021, mappando i metodi e le tecniche adottati dalla comunità scientifica per affrontare le problematiche specifiche della rilevazione di anomalie su dati generati da dispositivi IoT. La survey esamina in particolare tre aspetti critici: la riduzione della dimensionalità dei dati sensoristici, la localizzazione delle anomalie all'interno delle serie temporali e le sfide legate al monitoraggio in tempo reale sotto vincoli di risorse computazionali e temporali. Gli autori classificano gli approcci in base alla tipologia di algoritmo impiegato --- metodi statistici, tecniche di machine learning tradizionale e architetture di deep learning --- evidenziando come gli approcci basati su autoencoder e reti ricorrenti (LSTM) abbiano mostrato risultati particolarmente promettenti per la rilevazione di anomalie su serie temporali multivariate.

\section{Sintesi e posizionamento della soluzione proposta}

Dall'analisi dello stato dell'arte emerge che le soluzioni esistenti affrontano in modo separato i tre ambiti coinvolti nel presente lavoro: il matching tra domanda e offerta di spedizioni, il monitoraggio IoT della catena del freddo e la rilevazione di anomalie basata su deep learning. Le piattaforme di trasporto collaborativo come Flock Freight e Quicargo offrono meccanismi di consolidamento dei carichi ma non integrano funzionalità di monitoraggio per spedizioni refrigerate. Il sistema IoT di Gillespie et al. dimostra l'efficacia del monitoraggio in tempo reale ma si limita a due sonde di temperatura con soglie statiche, senza sfruttare le potenzialità del machine learning per la rilevazione di anomalie multivariate.

La review sistematica di Sgueglia et al. conferma come gli approcci basati su LSTM Autoencoder rappresentino attualmente una delle tecniche più promettenti per l’anomaly detection in serie temporali IoT multivariate, evidenziando tuttavia che la maggior parte dei lavori considerati propone componenti analitici isolati, senza integrazione in piattaforme operative complete \cite{sgueglia2022}.

ChillChain si propone di colmare questo gap attraverso una piattaforma integrata che combina un marketplace con pricing dinamico basato sull'occupazione del veicolo, un sistema di monitoraggio in tempo reale e un servizio di rilevazione anomalie basato su LSTM Autoencoder.